\section{Функция GetLoadFileNameExplorer}
\footnotesize
\begin{listing}{1}
/*************************************************************************
* Функция       : GetLoadFileNameExplorer                                *
*------------------------------------------------------------------------*
* Назначение    : Диалог выбора файла для открытия                       *
* Входные       : hWnd - дескриптор окна, outPath - путь к файлу         *
* Выходные      : bool - успех операции                                  *
*------------------------------------------------------------------------*
* ОПИСАНИЕ ФУНКЦИОНИРОВАНИЯ:                                             *
* Обертка над Win32 API для вызова стандартного диалога открытия файла.  *
*************************************************************************/
bool GetLoadFileNameExplorer(HWND hWnd, wstring& outPath) {
    wchar_t filename[MAX_PATH] = L""; 
    OPENFILENAMEW ofn; 
    ZeroMemory(&ofn, sizeof(ofn));
    ofn.lStructSize = sizeof(ofn); 
    ofn.hwndOwner = hWnd;
    ofn.lpstrFilter = L"Текстовые файлы (*.txt)\0*.txt\0Все файлы (*.*)\0*.*\0";
    ofn.lpstrFile = filename; 
    ofn.nMaxFile = MAX_PATH;
    ofn.Flags = OFN_FILEMUSTEXIST | OFN_PATHMUSTEXIST;
    if (GetOpenFileNameW(&ofn)) { 
        outPath = filename; 
        return true; 
    }
    return false;
}
\end{listing}
\normalsize

\pagebreak

\section{Функция GetSaveFileNameExplorer}
\footnotesize
\begin{listing}{1}
/*************************************************************************
* Функция       : GetSaveFileNameExplorer                                *
*------------------------------------------------------------------------*
* Назначение    : Диалог выбора файла для сохранения                     *
* Входные       : hWnd - дескриптор окна, outPath - путь к файлу         *
* Выходные      : bool - успех операции                                  *
*------------------------------------------------------------------------*
* ОПИСАНИЕ ФУНКЦИОНИРОВАНИЯ:                                             *
* Обертка над Win32 API для вызова стандартного диалога сохранения файла. *
*************************************************************************/
bool GetSaveFileNameExplorer(HWND hWnd, wstring& outPath) {
    wchar_t filename[MAX_PATH] = L"output.txt"; 
    OPENFILENAMEW ofn; 
    ZeroMemory(&ofn, sizeof(ofn));
    ofn.lStructSize = sizeof(ofn); 
    ofn.hwndOwner = hWnd;
    ofn.lpstrFilter = L"Текстовые файлы (*.txt)\0*.txt\0";
    ofn.lpstrFile = filename; 
    ofn.nMaxFile = MAX_PATH;
    ofn.Flags = OFN_OVERWRITEPROMPT | OFN_PATHMUSTEXIST;
    if (GetSaveFileNameW(&ofn)) { 
        outPath = filename; 
        return true; 
    }
    return false;
}
\end{listing}
\normalsize

\pagebreak

\section{Функция ImportDataFromFile}
\footnotesize
\begin{listing}{1}
/*************************************************************************
* Функция       : ImportDataFromFile                                     *
*------------------------------------------------------------------------*
* Назначение    : Импорт и парсинг данных из текстового файла            *
* Входные       : filename - путь к файлу                                *
* Выходные      : bool - успех импорта                                   *
*------------------------------------------------------------------------*
* ОПИСАНИЕ ФУНКЦИОНИРОВАНИЯ:                                             *
* Читает вариант задания и формулы полей, выполняет их первичный парсинг. *
*************************************************************************/
bool ImportDataFromFile(const wstring& filename) {
    ifstream file(filename);
    if (!file.is_open()) return false;
    string line;
    if (getline(file, line)) { 
        try { g_VariantNum = stoi(line); } 
        catch (...) { return false; } 
    }
    if (getline(file, line)) 
        g_FieldAStr = wstring(line.begin(), line.end());
    if (getline(file, line)) 
        g_FieldBStr = wstring(line.begin(), line.end());
    file.close();
    if (g_VariantNum <= 0 || g_FieldAStr.empty() || g_FieldBStr.empty()) 
        return false;
    FieldMathCalculations::ParsePolynomials(g_FieldAStr, g_FieldBStr);
    return true;
}
\end{listing}
\normalsize

\pagebreak

\section{Функция ExportDataToFile}
\footnotesize
\begin{listing}{1}
/*************************************************************************
* Функция       : ExportDataToFile                                       *
*------------------------------------------------------------------------*
* Назначение    : Экспорт отчета в файл (UTF-8 с BOM)                    *
* Входные       : filename - путь к файлу                                *
* Выходные      : bool - успех записи                                    *
*------------------------------------------------------------------------*
* ОПИСАНИЕ ФУНКЦИОНИРОВАНИЯ:                                             *
* Конвертирует отчет в UTF-8 и записывает в файл с добавлением BOM-метки. *
*************************************************************************/
bool ExportDataToFile(const wstring& filename) {
    wstring report = BuildReportText();
    if (report.empty() || g_VariantNum == 0) return false;
    int s_n = WideCharToMultiByte(CP_UTF8, 0, &report[0], 
                                  (int)report.size(), NULL, 0, NULL, NULL);
    string utf8_str(s_n, 0);
    WideCharToMultiByte(CP_UTF8, 0, &report[0], (int)report.size(), 
                        &utf8_str[0], s_n, NULL, NULL);                  
    ofstream file(filename, ios::binary);
    if (!file.is_open()) return false;
    const unsigned char bom[] = { 0xEF, 0xBB, 0xBF };
    file.write((const char*)bom, sizeof(bom)); 
    file.write(utf8_str.c_str(), utf8_str.size()); 
    file.close();
    return true;
}
\end{listing}
\normalsize

\pagebreak

\section{Функция BuildInputDataText}
\footnotesize
\begin{listing}{1}
/*************************************************************************
* Функция       : BuildInputDataText                                     *
*------------------------------------------------------------------------*
* Назначение    : Формирование текстового блока исходных данных          *
* Входные       : Нет                                                    *
* Выходные      : wstring - отформатированный текст                      *
*------------------------------------------------------------------------*
* ОПИСАНИЕ ФУНКЦИОНИРОВАНИЯ:                                             *
* Создает текстовое представление входных данных для отображения в GUI.  *
*************************************************************************/
wstring BuildInputDataText() {
    wstringstream ss;
    ss << L"====================================================================\n";
    ss << L"        ИСХОДНЫЕ ДАННЫЕ ВАРИАНТА\n";
    ss << L"====================================================================\n\n";
    ss << L"  ► Загружен вариант №:  " << g_VariantNum << L"\n\n";
    ss << L"  ► Векторное поле ā (исследуется на поток и циркуляцию):\n";
    ss << L"    ā = " << g_FieldAStr << L"\n\n";
    ss << L"  ► Векторное поле b̄ (исследуется на потенциальность):\n";
    ss << L"    b̄ = " << g_FieldBStr << L"\n\n";
    ss << L"--------------------------------------------------------------------\n";
    ss << L"  Всё готово! Нажмите кнопку «Расчёт данных» для получения отчета.";
    return ss.str();
}
\end{listing}
\normalsize

\pagebreak